\documentclass[11pt, a4paper]{article}

% --- UNIVERSAL PREAMBLE BLOCK ---
\usepackage[a4paper, top=2.5cm, bottom=2.5cm, left=2cm, right=2cm]{geometry}
%%\usepackage{fontspec}
%%\usepackage[english, bidi=basic, provide=*]{babel}

%%\babelprovide[import, onchar=ids fonts]{english}

% Set default/Latin font to Sans Serif for technical clarity
%%\babelfont{rm}{Noto Sans}

\usepackage{amsmath}
\usepackage{amssymb}
\usepackage{booktabs}
\usepackage{listings}
\usepackage{xcolor}

% Code listing configuration
\definecolor{codegray}{rgb}{0.5,0.5,0.5}
\definecolor{codeblue}{rgb}{0.1,0.2,0.6}
\lstset{
  language=C++,
  basicstyle=\small\ttfamily,
  keywordstyle=\color{codeblue}\bfseries,
  commentstyle=\color{codegray}\itshape,
  breaklines=true,
  tabsize=2,
  frame=single,
  captionpos=b
}

\title{Line Intersection and Least Squares Approximation}
\author{Technical Brief}
\date{\today}

\begin{document}

\maketitle

\section{Intersection of Two Lines}
In a deterministic 2D setting, two lines $L_1$ and $L_2$ defined by the equations $a_i x + b_i y = c_i$ intersect at a unique point $\mathbf{x}$ if the system is non-singular. 

In higher dimensions or parametric contexts, we represent lines as:
\begin{equation}
    \mathbf{L}_i(t) = \mathbf{p}_i + t\mathbf{d}_i
\end{equation}
where $\mathbf{p}_i$ is a point on the line and $\mathbf{d}_i$ is a unit direction vector ($\|\mathbf{d}_i\| = 1$). For two lines in 3D, a unique intersection exists only if the lines are coplanar and non-parallel. If they are skew, we often seek the point of closest approach.

\section{Least Squares for Multiple Lines}
In practical applications such as computer vision (triangulation) or robotics, we often have $N$ lines that, due to noise, do not intersect at a single point. We seek a point $\mathbf{x}$ that minimizes the sum of squared distances to all lines.

\subsection{Objective Function}
The vector from a point $\mathbf{p}_i$ on the line to the target point $\mathbf{x}$ is $(\mathbf{x} - \mathbf{p}_i)$. The projection of this vector onto the line's direction $\mathbf{d}_i$ is:
\begin{equation}
    \text{proj}_{\mathbf{d}_i}(\mathbf{x} - \mathbf{p}_i) = ((\mathbf{x} - \mathbf{p}_i) \cdot \mathbf{d}_i)\mathbf{d}_i
\end{equation}
The squared perpendicular distance $D_i^2$ from $\mathbf{x}$ to the $i$-th line is the magnitude of the rejected component:
\begin{equation}
    D_i^2 = \|(\mathbf{x} - \mathbf{p}_i) - ((\mathbf{x} - \mathbf{p}_i) \cdot \mathbf{d}_i)\mathbf{d}_i\|^2
\end{equation}
Using the identity $(\mathbf{x} \cdot \mathbf{d})\mathbf{d} = (\mathbf{d}\mathbf{d}^T)\mathbf{x}$, we define the projection matrix $P_i = \mathbf{d}_i\mathbf{d}_i^T$. The distance squared can be rewritten as:
\begin{equation}
    D_i^2 = \|(I - P_i)(\mathbf{x} - \mathbf{p}_i)\|^2
\end{equation}

\subsection{Normal Equations}
To find the global minimum, we minimize $S(\mathbf{x}) = \sum_{i=1}^N D_i^2$. Since $(I - P_i)$ is symmetric and idempotent ($(I-P_i)^2 = I-P_i$), the gradient $\nabla S$ with respect to $\mathbf{x}$ is:
\begin{equation}
    \nabla S = 2 \sum_{i=1}^N (I - P_i)(\mathbf{x} - \mathbf{p}_i) = 0
\end{equation}
Rearranging into the form $A\mathbf{x} = \mathbf{b}$ yields:
\begin{equation}
    \left( \sum_{i=1}^N (I - \mathbf{d}_i\mathbf{d}_i^T) \right) \mathbf{x} = \sum_{i=1}^N (I - \mathbf{d}_i\mathbf{d}_i^T) \mathbf{p}_i
\end{equation}
This is a $3 \times 3$ (or $2 \times 2$) linear system that can be solved efficiently using Cholesky (LLT) or LDLT decomposition.

\section{Modern C++ Implementation}
The following implementation utilizes \texttt{std::transform\_reduce} with \texttt{std::execution::par\_unseq} for parallelized construction of the system matrices, reflecting modern C++20 paradigms.

\begin{lstlisting}
#include <vector>
#include <numeric>
#include <execution>
#include <Eigen/Dense>

struct Line {
  Eigen::Vector3d p; // Point on line
  Eigen::Vector3d d; // Unit direction vector
};

Eigen::Vector3d findCommonIntersection(const std::vector<Line>& lines) {
  struct Accumulator {
    Eigen::Matrix3d m_sum = Eigen::Matrix3d::Zero();
    Eigen::Vector3d v_sum = Eigen::Vector3d::Zero();
  };

  auto result = std::transform_reduce(
    std::execution::par_unseq,
    lines.begin(), lines.end(),
    Accumulator{},
    [](Accumulator a, Accumulator b) {
      return Accumulator{a.m_sum + b.m_sum, a.v_sum + b.v_sum};
    },
    [](const Line& l) {
      Eigen::Matrix3d I_minus_ddT = Eigen::Matrix3d::Identity() - l.d * l.d.transpose();
      return Accumulator{I_minus_ddT, I_minus_ddT * l.p};
    }
  );

  return result.m_sum.ldlt().solve(result.v_sum);
}
\end{lstlisting}

\end{document}
